\scanpage{105}
\mspage{94}

be a basis of the group $H^0(V,O_V[N\widetilde H])$. We have:
\[
f_*(O_V[N\widetilde H])\simeq(f_*O_V)[NH]_{\bar V}.
\]

From the isomorphism of ringed spaces $(V,f^*O_{\bar V})$ and $(\bar V,O_{\bar V})$ we obtain the isomorphism of the sheaves $O_V[N\widetilde H]$ and $(f_*O_V)[NH]_{\bar V}$.

Let $\bar\alpha,\bar y,\bar y_1,\ldots,\bar y_M$ be the elements over $\bar V$, corresponding to the elements $\alpha,y,y_1,\ldots,y_M$. Since $\alpha y\in O_V[N\widetilde H](P)$, then
\[
\bar\alpha\bar y\in(f_*O_V)[NH]_{\bar V}(\bar P).
\]

According to the definition of the number $N$, there exist $\beta_\nu\in O_{\bar V,\bar P}$ such that
\[
\bar\alpha\bar y=\sum_{\nu=1}^M\beta_\nu\bar y_\nu
\]

Hence, if we denote by the same symbol the functions, corresponding to $\beta_\nu$ at the mapping $f$, we obtain:
\[
\alpha y=\sum_{\nu=1}^M\beta_\nu y_\nu
\]
or, since $y\ne0$ in $P$
\[
\alpha=\sum_{\nu=1}^M\beta_\nu\frac{y_\nu}{y}.\tag{24}
\]

\margin{104}
